% paper1_quantum_walks_and_graph_laws.tex
% REVTeX 4-2, Physical Review A style.
% Scope: Formulas 3-21, 24, 32, 37-40 of formulas.txt (Formula 2 deliberately excluded).
\documentclass[aps,pra,10pt,twocolumn,superscriptaddress,nofootinbib]{revtex4-2}

\usepackage{amsmath}
\usepackage{amssymb}
\usepackage{booktabs}
\usepackage[colorlinks=true,allcolors=blue]{hyperref}

\newcommand{\E}{\mathbb{E}}
\newcommand{\pbar}{\bar{p}}
\newcommand{\Cl}{\operatorname{Cl}_2}
\newcommand{\erfc}{\operatorname{erfc}}
\newcommand{\sq}{\sqrt{3}}
\newcommand{\YQ}{Y^{Q}}

\begin{document}

\title{Exact localization of Grover quantum walks on crystal lattices:\\
transfer currents, a universal bound, and uniform-spanning-tree degree laws}

\author{Jacob Goodchild}
\affiliation{Independent researcher}

\date{September 28, 2026}

\begin{abstract}
Discrete-time Grover walks on periodic graphs localize: a finite fraction of the
wave function stays at the starting vertex forever because the walk operator
has flat (momentum-independent) eigenvalues. We compute the long-time average
return probability $\pbar$ exactly for a large family of lattices. For the
moving-shift walk on $\mathbb{Z}^d$ we show that
$\pbar$ is governed by $A_d=\int_0^\infty \erfc(\sqrt{s})^d\,ds$, and we obtain
$A_3=\tfrac12-(3-\sq)/\pi$,
$A_4=\tfrac12-(18-8\sq)/(3\pi)$ and an arccosine closed form for $A_5$.
For the flip-flop walk we prove that the $\pm1$ flat-band projectors are
$I-Y$ and $I-\YQ$, where $Y$ and $\YQ$ are the Kirchhoff transfer-current matrices
of the Laplacian and signless Laplacian. This reduces the trapping probability to
lattice Green functions. It gives closed forms in terms of Watson integrals
(simple cubic, body-centred cubic, face-centred cubic), complete elliptic
integrals (triangular, kagome, checkerboard, star, line graphs), and elementary
functions of the bond anisotropy. We prove the universal lower bound
$\pbar\ge (d-2)^2/[2d(d-1)]$ for edge-transitive $d$-regular lattices, where $d$
is the vertex degree. The bound is attained exactly on bipartite lattices whose vertex
stabilizer acts 2-transitively on the incident edges (honeycomb, diamond, the
Laves graph, $d$-dimensional diamond). We also solve lazy (Szegedy) versions on the
square, triangular, honeycomb and diamond lattices as explicit functions of the
laziness. The same transfer-current matrix yields the vertex-degree distribution
of the uniform spanning tree. On 2-arc-transitive lattices it is exactly
$1+\mathrm{Binomial}(z-1,1/(z-1))$, and we give closed forms for the simple-cubic,
BCC, FCC, pyrochlore, king and snub-square lattices. All closed forms were
identified by integer-relation (PSLQ) searches at 25--80 digits and checked
independently by Brillouin-zone diagonalization, direct simulation, or both.
\end{abstract}

\maketitle

%=====================================================================
\section{Introduction}
\label{sec:intro}
%=====================================================================

Quantum walks are the unitary analogues of classical random
walks~\cite{Kempe2003}. They underpin quantum search algorithms, in particular
Szegedy's quantization of Markov chains~\cite{Szegedy2004} and the lackadaisical
(lazy) walks introduced by Wong~\cite{Wong2015}. A striking feature of the Grover
walk in two or more dimensions is \emph{localization}: even on a perfect lattice,
the probability of finding the walker at its starting site does not decay to zero.
Inui, Konishi and Konno established this for the two-dimensional Grover walk
and computed the limiting return probability~\cite{Inui2004}. Higuchi, Konno, Sato
and Segawa related localization of Grover walks on crystal lattices to flat
bands of the walk operator and gave integral descriptions in any
dimension~\cite{Higuchi2014}. What has been missing is a systematic set of
\emph{closed forms} beyond the square lattice.

This paper supplies them. The main structural observation (Sec.~\ref{sec:model}) is
that for the flip-flop Grover walk the $+1$ and $-1$ flat-band projectors are the
projectors onto the cycle spaces of the Laplacian $L=D-A$ and of the signless
Laplacian $Q=D+A$. Their matrix elements between arcs at a vertex are therefore
Kirchhoff transfer currents~\cite{Kirchhoff1847,BurtonPemantle1993}, the same
objects that give edge correlations of the uniform spanning tree (UST). Three
consequences follow.
(i) Exact trapping probabilities reduce to lattice Green functions at the two
band edges. These are classical in several cases: Watson
integrals~\cite{Watson1939,GlasserZucker1977}, Joyce's cubic Green
functions~\cite{Joyce1994,Joyce1998} and Horiguchi's triangular and honeycomb
formulas~\cite{Horiguchi1972}.
(ii) Foster's theorem~\cite{Foster1949} and the Cauchy--Schwarz inequality give a
universal lower bound, attained on 2-arc-transitive bipartite lattices
(Sec.~\ref{sec:bound}).
(iii) The UST degree distribution at a vertex is
$\det[I+(x-1)Y]$~\cite{BurtonPemantle1993}, which gives an exact binomial law in the
symmetric case (Sec.~\ref{sec:ust}).

We also treat the moving-shift walk on $\mathbb{Z}^d$ in any dimension
(Sec.~\ref{sec:higherd}). There the trapping integral is a Cauchy or Gaussian average that
closes in elementary terms for $d\le5$.

\emph{Status of novelty.} Every result below was checked against the literature by
web search. Items we believe are already known, or probably known, are flagged
explicitly: the two-dimensional moving-shift values~\cite{Inui2004}, the integral
$A_5$ (which is equivalent to a mean-square maximum of half-normal variables
treated by Finch~\cite{Finch2011}), Kenyon's isoradial diagonal transfer
currents~\cite{Kenyon2002}, and the mean-field limit of the UST degree law. For the
remaining statements our searches found no prior closed form. Novelty should be
understood in that limited sense.

%=====================================================================
\section{Model setup}
\label{sec:model}
%=====================================================================

\subsection{Grover walks on periodic graphs}

Let $G$ be a locally finite periodic graph. The walk lives on arcs (directed edges)
$a=(o\to h_a)$, and each vertex $v$ of degree $d_v$ carries the Grover coin
$C_v=\tfrac{2}{d_v}J-I$. In the \emph{flip-flop} walk, the shift sends the arc
$(o\to h)$ to its reverse $(h\to o)$. In the \emph{moving-shift} walk on
$\mathbb{Z}^d$, the coin state $\pm e_i$ is transported by $\pm e_i$ and keeps its label.
For an initial state $\psi_0$ localized at $o$, the long-time average return
probability is
\begin{equation}
\pbar=\lim_{T\to\infty}\frac1T\sum_{t=0}^{T-1}\sum_{a:\,\text{tail}(a)=o}
|\langle a|U^t\psi_0\rangle|^2 .
\label{eq:pbar-def}
\end{equation}
After Fourier transformation, the continuous spectrum of $U(k)$ contributes
nothing to Eq.~\eqref{eq:pbar-def} by the Riemann--Lebesgue lemma. Only the flat
eigenvalues $\lambda$ (eigenvalues of $U(k)$ that do not depend on $k$) survive, and
\begin{equation}
\pbar=\sum_{\lambda\ \text{flat}}\ \sum_{a\ \text{at}\ o}
\bigl|\langle a|\,\overline{\Pi}_\lambda\,|\psi_0\rangle\bigr|^2,
\qquad
\overline{\Pi}_\lambda=\E_k\,\Pi_\lambda(k),
\label{eq:pbar-flat}
\end{equation}
where $\E_k$ denotes the normalized Brillouin-zone average.

\subsection{Flat bands as cycle spaces: the transfer-current principle}

\emph{Theorem 1 (transfer-current form of flip-flop trapping).}
For the flip-flop Grover walk on any periodic graph, starting on arc $a$ at vertex $o$,
\begin{align}
\pbar(a)&=\tfrac14\sum_{b\ \text{at}\ o}(\delta_{ab}-Y_{ab})^2
+\tfrac14\sum_{b\ \text{at}\ o}(\delta_{ab}-\YQ_{ab})^2 \nonumber\\
&\quad+\sum_{\mu}\,(\text{flat-band terms}\ \ge 0).
\label{eq:principle}
\end{align}
Here $Y_{ab}$ is the current through arc $b$ when a unit current enters at $o$
and leaves at $h_a$ in a unit-resistor network (the Kirchhoff transfer
current). $\YQ$ is the same object built from $Q=D+A$ in place of $L=D-A$. The
last sum runs over flat bands $\mu$ of the simple random walk $D^{-1/2}AD^{-1/2}$.

\emph{Derivation outline.}
(1) Let $u$ be the vector with components $d_v^{-1/2}$ on the arcs at $v$, and let $S$ be the
flip-flop shift, so $U=S(2uu^{*}-I)$. The flat eigenvalue $+1$ is supported on
$K=\{\psi:\langle u,\psi\rangle=0=\langle u,S\psi\rangle\}$, where $U=-S$.
(2) Divergence-free (circulating) flows satisfy both constraints. The $+1$ eigenspace
is therefore the cycle space of the graph, and its projector is
$\tfrac12(I-S)-\tfrac12(I-S)uu^{*}(I-S)/\langle u,(I-S)u\rangle$.
(3) Averaged over $k$, the projector onto the orthogonal complement of the cycle space
(the cut space) has matrix elements $Y_{ab}=\tfrac12[R(o,h_a)+R(o,h_b)-R(h_a,h_b)]$,
written with the resistance distance $R$, that is, with the Green function
$G(x)=\E_k[\cos(k\cdot x)/\sum_r(1-\cos k\cdot\delta_r)]$.
(4) Replacing $S$ by $-S$ interchanges $L$ and $Q$, which gives the $-1$ channel.
For bipartite graphs $Q$ is unitarily equivalent to $L$, so the two channels are equal.
(5) Every other flat eigenvalue of $U$ comes, through the Szegedy correspondence,
from a flat band $\mu$ of the random walk, with $\lambda=\mu\pm i\sqrt{1-\mu^2}$.

Two identities hold throughout:
\begin{align}
\sum_{b}Y_{ab}&=\sum_b\YQ_{ab}=1, \label{eq:rowsum}\\
\sum_{e\in\text{cell}}Y_{ee}&=\sum_{e\in\text{cell}}\YQ_{ee}=|V_{\text{cell}}|.
\label{eq:foster}
\end{align}
Equation~\eqref{eq:rowsum} is current conservation. Equation~\eqref{eq:foster} is Foster's
theorem~\cite{Foster1949} per unit cell (the trace of $L^{-1}L$). On an edge-transitive
$d$-regular lattice it forces $Y_{aa}=\YQ_{aa}=2/d$.

\subsection{Flip-flop walks on the cubic lattices}

For a Bravais lattice with neighbour vectors $\pm\delta_1,\dots,\pm\delta_P$,
Eq.~\eqref{eq:principle} becomes $\pbar(\phi)=\|M_+\phi\|^2+\|M_-\phi\|^2$ with
$M_+=\tfrac12 I-N$ and
\begin{equation}
N_{(s,p),(t,q)}=\tfrac14\bigl[G(0)-G(\delta_p)-G(\delta_q)+G(t\delta_q-s\delta_p)\bigr].
\end{equation}

\emph{Simple cubic.} Let
\begin{equation}
W=\frac{\sqrt6}{32\pi^3}\,\Gamma(\tfrac1{24})\Gamma(\tfrac5{24})\Gamma(\tfrac7{24})\Gamma(\tfrac{11}{24})
=1.5163860591519780\ldots
\end{equation}
be Watson's simple-cubic integral~\cite{Watson1939,GlasserZucker1977}. Then
$N_{\rm diag}=1/6$, $N(+i,-i)=\alpha=(7W+54/(\pi^2W)-12)/36$, and
$N(i,j\ne i)=\beta=(24-7W-54/(\pi^2W))/144$, with $\alpha+4\beta=1/3$. Starting in one coin
state,
\begin{equation}
\pbar_{\rm SC}=2\Bigl[\tfrac19+\alpha^2+4\beta^2\Bigr]=0.26672720352085343375\ldots
\label{eq:sc}
\end{equation}
The uniform coin state escapes completely ($\pbar=0$).

\emph{Body-centred cubic.} With $W_b=\Gamma(\tfrac14)^4/(4\pi^3)=1.39320392968567685918\ldots$,
we have $N_{\rm diag}=1/8$ and
$\alpha=(3W_b+36/(\pi^2W_b)-6)/16$,
$\beta_s=(2-W_b+4/(\pi^2W_b))/16$ and
$\beta_d=(1-8/(\pi^2W_b))/8$. Hence
\begin{equation}
\pbar_{\rm BCC}=2\Bigl[\tfrac9{64}+\alpha^2+3\beta_s^2+3\beta_d^2\Bigr]
=0.32150479124143771009\ldots
\end{equation}

\emph{Face-centred cubic.} FCC is not bipartite, so the $-1$ channel requires the
Green function at the opposite band edge, $G'(x)=\E_k[\cos(k\cdot x)/(6+2s)]$ with
$s=c_1c_2+c_2c_3+c_3c_1$ and $c_i=\cos k_i$. We find that the far-edge value is governed by
the \emph{simple-cubic} Watson constant:
\begin{equation}
G'(0)=\frac{W}{6\sqrt2}
=\frac{\sq\,\Gamma(\tfrac1{24})\Gamma(\tfrac5{24})\Gamma(\tfrac7{24})\Gamma(\tfrac{11}{24})}{192\pi^3}.
\label{eq:fcc-far}
\end{equation}
This identity probably follows from Joyce's analysis of the FCC Green
function~\cite{Joyce1998}. With $Y=\sqrt2W$ we further have $G'(1,1,0)=(2-Y)/12$,
$G'(2,0,0)=(7Y+108/(\pi^2Y)-36\sq/\pi)/72$,
$G'(2,1,1)=(2Y-3-27/(\pi^2Y))/9$, and
$G'(2,2,0)=(Y+54/(\pi^2Y)+6\sq/\pi-8)/6$. All of these satisfy the FCC lattice equations to a residual of
$3\times10^{-42}$. With
\begin{equation}
W_f=\frac{3\,\Gamma(\tfrac13)^6}{2^{14/3}\pi^4}=0.44822039439\ldots,
\end{equation}
the ordinary channel follows from $W_f$ and $1/(\pi^2W_f)$. The totals are
\begin{align}
\pbar_{\rm FCC}&=0.18948881307124000416+0.18942581798307167926\nonumber\\
&=0.37891463105431168342\ldots
\end{align}

\subsection{Two-dimensional lattices}

\emph{Triangular.} The Laplacian resistances are
$a(1)=1/3$, $a(\sq)=2\sq/\pi-2/3$ and $a(2)=8/3-4\sq/\pi$. At the far band edge we find
\begin{equation}
X\equiv\E_k\Bigl[\frac1{3+\cos k_1+\cos k_2+\cos(k_1-k_2)}\Bigr]
=\frac{2^{2/3}\Gamma(\tfrac13)^3}{8\pi^2},
\end{equation}
with $G'(1)=(1-3X)/3$, $G'(\sq)=(9X-2-\sq/(\pi X))/3$ and $G'(2)=(3X-4+2\sq/(\pi X))/3$. The two channels are
\begin{align}
\text{(+1)}&=\frac{43\pi^2-120\sq\pi+324}{72\pi^2},\\
\text{($-$1)}&=X^2-\frac X2+\frac{19}{72}-\frac{\sq}{6\pi}-\frac{\sq}{18\pi X}+\frac1{8\pi^2X^2},
\end{align}
with total $\pbar=0.26795570253708055129\ldots$.
The square lattice gives $\pbar=\tfrac12-\tfrac2\pi+\tfrac3{\pi^2}$, which is very likely known.

\emph{Honeycomb and kagome.} On the honeycomb lattice, Eq.~\eqref{eq:foster} and edge
transitivity force $Y=(2/3,1/6,1/6)$, so $\pbar=1/12$ exactly. Kagome has a flat band
$\mu=-1/2$ of the random walk. This produces two extra walk eigenvalues
$\lambda=e^{\pm2\pi i/3}$, each contributing
$(19-42\sq/\pi+108/\pi^2)/648$. With $K=K(m)$ and $E=E(m)$ at parameter
$m=5/32$ (we use the parameter convention $K(m)=\int_0^{\pi/2}(1-m\sin^2\theta)^{-1/2}d\theta$
throughout), the total is
\begin{align}
\pbar_{\rm kag}&=\frac{103}{144}-\frac{7\sq}{48\pi}+\frac3{8\pi^2}
-\frac{23\sqrt2E}{12\pi}+\frac{127\sqrt2K}{288\pi}\nonumber\\
&\quad+\frac{4E^2}{\pi^2}-\frac{13KE}{6\pi^2}+\frac{61K^2}{144\pi^2}\nonumber\\
&=0.19146243000085471390\ldots
\end{align}
The elliptic input is the identity
$\E_k[1/(11-\cos k_1-\cos k_2-\cos(k_1-k_2))]=K(5/32)/(4\sqrt2\pi)$.

\emph{Lieb, dice, 4.8.8, checkerboard and star.} Table~\ref{tab:2d} collects
these results. We write $c=\sqrt2\arccos(1/3)/\pi$. The 4.8.8 Laplacian determinant
factorizes as $64-4(3+\cos k_1)(3+\cos k_2)$, which is the origin of $\arccos(1/3)$. For the star
lattice, $\kappa=\sqrt6\,K(3/128)/\pi$.

\begin{table*}
\caption{Exact long-time return probabilities of the flip-flop Grover walk on
two-dimensional lattices, from Eq.~\eqref{eq:principle}. $c=\sqrt2\arccos(1/3)/\pi$,
$\kappa=\sqrt6K(3/128)/\pi$, and $K,E$ in the checkerboard rows are at parameter $m=1/4$.}
\label{tab:2d}
\begin{ruledtabular}
\begin{tabular}{lll}
Lattice, start & Closed form & Value\\
\hline
Lieb, corner & $(1-4/\pi+6/\pi^2)/4$ & $0.0836718892797159856$\\
Lieb, edge site & $1/8$ & $0.125$\\
Dice, hub & $(31-20\sq/\pi+45/\pi^2)/144$ & $0.17036718996853721$\\
Dice, rim & $11/48$ & $0.2291\overline{6}$\\
4.8.8, square edge & $(4-6c+3c^2)/16$ & $0.0997756085862333247$\\
4.8.8, bridge edge & $3c^2/16$ & $0.0575730175785728199$\\
Star 3.12.12, triangle edge & $1883/10800-7\kappa/96+3\kappa^2/128$ & $0.1200922563314837021$\\
Star 3.12.12, bridge edge & $107/2700+\kappa/48+3\kappa^2/128$ & $0.1008716548141747862$\\
Checkerboard, collinear &
$\frac{2097}{1024}-\frac{185}{128\pi}+\frac{79}{32\pi^2}+\frac{123K}{32\pi}-\frac{93E}{8\pi}
+\frac{39K^2}{16\pi^2}-\frac{57KE}{4\pi^2}+\frac{43E^2}{2\pi^2}$ & $0.2915287743505391057$\\
Checkerboard, perpendicular &
$\frac{777}{1024}-\frac{37}{64\pi}+\frac{25}{32\pi^2}+\frac{45K}{64\pi}-\frac{33E}{16\pi}
+\frac{3K^2}{4\pi^2}-\frac{3KE}{\pi^2}+\frac{7E^2}{2\pi^2}$ & $0.2954779129769429679$\\
\end{tabular}
\end{ruledtabular}
\end{table*}

\subsection{Line graphs}

Let $L(G)$ be the line graph of a Bravais lattice $G$ with bond vectors $e_1,\dots,e_P$, and put
$B(k)=(1+e^{ik\cdot e_t})_{t=1}^P$. Then $A_{L(G)}+2=B^\dagger B$, the degree is $4P-2$, and
\begin{align}
L^{-1}&=\tfrac1{4P}\bigl[I+B^\dagger B/(4P-BB^\dagger)\bigr],\\
Q^{-1}&=\tfrac1{4P-4}\bigl[I-B^\dagger B/(4P-4+BB^\dagger)\bigr],\\
F&=I-B^\dagger B/(BB^\dagger),
\end{align}
where $F$ is the projector onto the flat band $\mu=-2/(4P-2)$. Every trapping channel on
$L(G)$ is therefore a finite combination of Green functions of $G$ at three energies.
The checkerboard ($=L(\mathbb{Z}^2)$) and kagome ($=L$(honeycomb)) are instances.

\emph{Line graph of $\mathbb{Z}^3$} (degree 10). Let $g(t)=\E[1/(t-\cos k_1-\cos k_2-\cos k_3)]$ and
$g_j=g^{(j)}(7)$. Every channel is a polynomial in $W$, $X=1/(\pi^2W)$ and the $g_j$. For the collinear start,
\begin{align}
(+1)&=\tfrac{65}{324}-\tfrac{109W}{2592}-\tfrac{9X}{16}+\tfrac{45701W^2}{2985984}
+\tfrac{9347WX}{27648}+\tfrac{1973X^2}{1024},\nonumber\\
\text{(flat)}&=\tfrac{259}{1296}-\tfrac{6233W}{41472}-\tfrac{453X}{256}
+\tfrac{773765W^2}{23887872}\nonumber\\
&\quad+\tfrac{166595WX}{221184}+\tfrac{37685X^2}{8192},
\end{align}
and $\pbar=0.38013642908365187096$ (perpendicular start: $0.38100581744260323594$).
Joyce's closed form for $g$~\cite{Joyce1994} makes the $-1$ channel a polynomial in
complete elliptic integrals.

\emph{Line graphs of BCC and FCC and of the triangular lattice.} These are solved completely
in the same way. The totals are listed in Table~\ref{tab:lg}. The BCC case needs, besides $W_b$,
the Green function $\E[1/(10+4c_1c_2c_3)]=\kappa^2/10$ with $\kappa=2K(m)/\pi$ and
$m=\tfrac12-\sqrt{21}/10$. The triangular case needs the Green function \emph{below} the band.
Its modulus is complex, but a modular transformation gives the real form
\begin{equation}
\E_k\Bigl[\frac1{7+\cos k_1+\cos k_2+\cos(k_1-k_2)}\Bigr]=\frac{K(m_*)}{135^{1/4}\pi},
\end{equation}
where $m_*=\tfrac12-\tfrac{23\sqrt{15}}{180}$ is a root of $2160m^2-2160m+11$.

\begin{table}
\caption{Flip-flop trapping on line graphs (start classes by the angle between the
start bond and the site's own bond).}
\label{tab:lg}
\begin{ruledtabular}
\begin{tabular}{llr}
Line graph & start & $\pbar$\\
\hline
$L(\text{triangular})$ & collinear & $0.37891749229289535007$\\
 & $60^\circ$ & $0.37983613893989158023$\\
 & $120^\circ$ & $0.38314651208017409854$\\
$L(\text{SC})$ & collinear & $0.38013642908365187096$\\
 & perpendicular & $0.38100581744260323594$\\
$L(\text{BCC})$ & collinear & $0.41704934756699764035$\\
 & $70.5^\circ$ & $0.41732104849596726107$\\
 & $109.5^\circ$ & $0.41775946732413853940$\\
$L(\text{FCC})$ & collinear & $0.44924601047168703813$\\
 & $60^\circ$ & $0.44932628436614800807$\\
 & $90^\circ$ & $0.44943990277365356359$\\
 & $120^\circ$ & $0.44977663121249284644$\\
\end{tabular}
\end{ruledtabular}
\end{table}

\subsection{Pyrochlore}

The pyrochlore lattice ($=L$(diamond), degree 6) has a doubly degenerate flat band
$\mu=-1/3$. Its projector is $F=I-B^\dagger(BB^\dagger)^{-1}B$, where $B$ is the diamond
incidence matrix. With $W=W_f$ and $X=1/(\pi^2W)$, the Laplacian transfer currents from an edge are
$1/3$ (self), $1/6$ (the two other edges of the same tetrahedron),
$(1+8W+3X)/48$ (collinear partner) and $(15-8W-3X)/96$ (the other two). The channels are
\begin{align}
(+1)&=\tfrac{W^2}{96}+\tfrac{WX}{128}-\tfrac{13W}{1152}+\tfrac{3X^2}{2048}-\tfrac{13X}{3072}+\tfrac{2531}{18432},\\
\text{(flat)}&=\tfrac{3W^2}{64}+\tfrac{9WX}{256}-\tfrac{13W}{256}+\tfrac{27X^2}{4096}-\tfrac{39X}{2048}+\tfrac{163}{4096}.
\end{align}
The $-1$ channel is a quadratic polynomial in $g_0,g_1,g_2$, the FCC Green function
$g(t)=\E[1/(t-s)]$ and its first two derivatives at $t=15$, which Joyce's formula makes explicit.
The total is $\pbar_{\rm pyro}=0.29388359496\ldots$.

\subsection{Anisotropic walks}

\emph{Square lattice.} Take the random walk with $P(\pm x)=\alpha/2$ and $P(\pm y)=\beta/2$
($\alpha+\beta=1$), its Szegedy walk, and $\theta=\arcsin\sqrt\alpha$. The key integral is
$\E[(1-\cos k_1)/(\alpha(1-\cos k_1)+\beta(1-\cos k_2))]=2\theta/(\pi\alpha)$. It gives
\begin{align}
Y(+x,+x)&=2\theta/\pi,\qquad
Y(+x,-x)=\tfrac2\pi(\cot\theta-\theta\cot^2\theta),\\
Y(+x,\pm y)&=\tfrac{\tan\theta}{2}\Bigl[1-\tfrac2\pi(\theta+\cot\theta-\theta\cot^2\theta)\Bigr],
\end{align}
and $\pbar(\alpha)=\tfrac12[(1-2\theta/\pi)^2+Y(+x,-x)^2+2Y(+x,y)^2]$. As checks,
$\alpha=1/4$ gives $0.25400036756659953284$ and $\alpha=0.7$ gives $0.10771461463761465554$.

\emph{Honeycomb.} Choose bond type $i$ with probability $p_i$, and set
$\theta_i=\arctan(\lambda p_i)$ with $\lambda=(p_1p_2p_3)^{-1/2}$, so that $\sum\theta_i=\pi$.
Write $t_i=\theta_i/\pi$. Then $T(1,1)=2t_1$ (Kenyon~\cite{Kenyon2002}) and
\begin{equation}
T(1,j)=\bigl[\tfrac12-p_k-p_1t_1-p_jt_j+p_kt_k\bigr]/p_1,\quad\{j,k\}=\{2,3\},
\end{equation}
so that $\pbar=\tfrac12[(1-2t_1)^2+(p_1/p_2)T(1,2)^2+(p_1/p_3)T(1,3)^2]$. For example,
$p=(1/2,1/4,1/4)$ gives $\pbar=0.04680597938234766442$.

\emph{Triangular.} For conductances $q_i$ with $\sum q_i=1$, set
$\theta_i=\arctan(\lambda q_i)$ with $\lambda=(q_1q_2+q_2q_3+q_3q_1)^{-1/2}$, so that $\sum\theta_i=\pi/2$.
Let $t_i=\theta_i/\pi$ and $c_1=\cot\theta_1$. The six currents at the origin for a unit
current from $o$ to its $+e_1$ neighbour are
\begin{align}
T(+e_1)&=2t_1,\qquad T(-e_1)=\tfrac2\pi(c_1-\theta_1c_1^2),\nonumber\\
T(+e_2)&=-\tfrac{q_2}{2q_3}+\tfrac{q_2(q_1+q_3)}{q_1q_3}t_1+\tfrac{q_2+q_3}{q_3}t_2,\nonumber\\
T(-e_2)&=\tfrac{q_2(q_1+q_3)}{q_1^2}t_1+\tfrac{q_2+q_3}{q_2}t_2-\tfrac{c_1}{\pi},\nonumber\\
T(+e_3)&=\tfrac12-\tfrac{q_1-q_3}{q_1}t_1-\tfrac{q_2+q_3}{q_2}t_2,
\end{align}
and $T(-e_3)=1-\sum(\text{others})$. The diagonal is Kenyon's isoradial theorem~\cite{Kenyon2002}.

\subsection{Lazy (Szegedy) flip-flop walks}
\label{sec:lazy}

Add a self-loop with holding probability $\varepsilon$ and use the Szegedy coin. Self-loops
carry no circulation, so the $+1$ channel is unchanged, while the $-1$ channel acquires a
loop amplitude.

\emph{Square lattice.} Let $k=(1-\varepsilon)/(1+\varepsilon)$, $\mathcal K=K(k^2)/\pi$ and
$\mathcal E=E(k^2)/\pi$. The amplitudes are
$A(+x)=\tfrac14+\tfrac12(1-k)\mathcal K$,
$A(-x)=-\tfrac14[1-2/k+2(k-1)\mathcal K+4\mathcal E/k]$,
$A(\pm y)=-\tfrac14[1+(2/k-2)\mathcal K-2\mathcal E/k]$ and
$A(\text{loop})=-\sqrt{(1-k)/(2k)}\,(\tfrac12-(1-k)\mathcal K)$. Then
\begin{equation}
\pbar(\varepsilon)=\tfrac14-\tfrac1\pi+\tfrac3{2\pi^2}+A(+x)^2+A(-x)^2+2A(y)^2+A(\text{loop})^2 .
\end{equation}

\emph{Triangular lattice.} Let $t=3(1+\varepsilon)/(1-\varepsilon)$ and $S=\cos k_1+\cos k_2+\cos(k_1-k_2)$.
For $t>3/2$ the Green function has the real-parameter form
\begin{equation}
g(t)=\E_k\Bigl[\frac1{t+S}\Bigr]=\frac{2K(m)}{\pi[(t-1)^3(t+3)]^{1/4}},
\end{equation}
\begin{equation}
m=\frac12-\frac{t^2-3}{2(t-1)^{3/2}\sqrt{t+3}}.
\label{eq:tri-real}
\end{equation}
We need $G_1=(1-tg)/3$, together with
\begin{align}
G_{\sq}&=-\tfrac23+\tfrac{2t^2+3t-6}{3}g+\tfrac{(t-1)(t+3)(2t-3)}{3}g',\\
G_2&=-\tfrac{2(t-1)}3-\tfrac{2t^2+4t-9}{3}g-\tfrac{2(t-1)(t+3)(2t-3)}{3}g'.
\end{align}
The $-1$ amplitudes are
$A(+e_1)=\tfrac13+(t-3)g/6$,
$A(-e_1)=-(g+2G_1+G_2)/4$,
$A(\text{adj})=-(g+3G_1)/4$ (two arcs),
$A(\text{rem})=-(g+2G_1+G_{\sq})/4$ (two arcs) and
$A(\text{loop})=-\tfrac{\sqrt{t-3}}6(1+(3-t)g)$, and
\begin{align}
\pbar(\varepsilon)&=\frac{43\pi^2-120\sq\pi+324}{72\pi^2}+A(+e_1)^2+A(-e_1)^2\nonumber\\
&\quad+2A(\text{adj})^2+2A(\text{rem})^2+A(\text{loop})^2 .
\end{align}
For example, $\pbar(0.2)=0.31653136185695229253$ and $\pbar(0.5)=0.35388644533661557374$.

\emph{Honeycomb.} Let $r=3(1+\varepsilon)/(1-\varepsilon)$, $m=16r/[(r-1)^3(r+3)]$ and
$g=4K(m)/(\pi\sqrt{(r-1)^3(r+3)})$. Then
\begin{align}
288\,\pbar(\varepsilon)&=g^2r^2(r-3)^2(r-1)(r+3)\nonumber\\
&\quad-4gr(r-4)(r-3)(r+3)+4(r^2-4r+9).
\end{align}
At $\varepsilon=0$ this gives $1/12$.

\emph{Diamond.} Let $r=4(1+\varepsilon)/(1-\varepsilon)$ and $\tau=(r^2-4)/4$, and let $g=\E[1/(\tau-s)]$ be the FCC Green
function in Joyce's closed form~\cite{Joyce1998},
$g=\tau^{-1}\mathcal L(3/\tau)$ with
$\mathcal L(x)=3^{3/2}(3+x)^{-3/2}(2-\sqrt{1-x})(2K(m)/\pi)^2$ and
$m=\tfrac12-\tfrac{\sq}{2}(3+x)^{-3/2}(4x+(3-x)\sqrt{1-x})$. Then
\begin{align}
3072\,\pbar(\varepsilon)&=g^2r^2(r-4)^2(r-2)(r+4)\nonumber\\
&\quad-8gr(r-6)(r-4)(r+4)+16(r^2-6r+40).
\end{align}
At $\varepsilon=0$ this gives $1/6$.

\emph{General $z$.} For a $z$-regular bipartite lattice of girth at least 6 whose vertex symmetry acts
2-transitively on the arcs, let $a_2(k)=|f(k)|^2-z$, $r=z(1+\varepsilon)/(1-\varepsilon)$,
$h=\E[1/(r^2-z-a_2)]$ and $u=hr(r-z)$. The amplitudes are
$A_{\rm start}=(2u+z-2)/(2z)$,
$A_{\rm other}=-[u(r-z+2)-r+2z-2]/[2z(z-1)]$ and
$A_{\rm loop}^2=(r-z)(1-u)^2/(2z^2)$, and
\begin{equation}
\pbar(\varepsilon)=\frac{(z-2)^2}{4z(z-1)}+A_{\rm start}^2+(z-1)A_{\rm other}^2+A_{\rm loop}^2 .
\end{equation}

%=====================================================================
\section{The universal localization bound theorem}
\label{sec:bound}
%=====================================================================

\emph{Theorem 2.} Consider the flip-flop Grover walk on a periodic graph, started on an arc $a$
at a vertex of degree $d$. Then
\begin{equation}
\pbar(a)\ \ge\ \frac{d}{4(d-1)}\Bigl[(1-Y_{aa})^2+(1-\YQ_{aa})^2\Bigr].
\label{eq:bound-general}
\end{equation}
If the lattice is edge-transitive and $d$-regular, where $d$ is the \emph{vertex degree}
and not the spatial dimension, then
\begin{equation}
\pbar\ \ge\ \frac{(d-2)^2}{2d(d-1)} .
\label{eq:bound}
\end{equation}

\emph{Proof.} The flat-band terms in Eq.~\eqref{eq:principle} are nonnegative, so it
suffices to bound each $\pm1$ channel. By Eq.~\eqref{eq:rowsum}, the $d-1$ off-diagonal
entries $Y_{ab}$ ($b\ne a$) sum to $1-Y_{aa}$. The Cauchy--Schwarz inequality then gives
$\sum_{b\ne a}Y_{ab}^2\ge(1-Y_{aa})^2/(d-1)$. Hence
\begin{align}
\tfrac14\sum_b(\delta_{ab}-Y_{ab})^2
&\ge\tfrac14\Bigl[(1-Y_{aa})^2+\tfrac{(1-Y_{aa})^2}{d-1}\Bigr]\nonumber\\
&=\tfrac{d}{4(d-1)}(1-Y_{aa})^2,
\end{align}
and likewise for $\YQ$. Under edge transitivity, Eq.~\eqref{eq:foster} gives $Y_{aa}=\YQ_{aa}=2/d$,
and substituting yields Eq.~\eqref{eq:bound}. $\square$

\emph{Equality.} Equality holds if and only if (i) all off-diagonal currents from $a$ are equal
in both channels and (ii) there are no additional flat bands. Both hold for bipartite
lattices whose vertex stabilizer permutes the $d$ arcs 2-transitively, because then $\YQ$
equals $Y$ up to signs. This gives the exact values
\begin{equation}
\begin{array}{ll}
\text{honeycomb }(d=3): & \pbar=1/12,\\
\text{Laves graph / srs }(d=3): & \pbar=1/12,\\
\text{diamond }(d=4): & \pbar=1/6,\\
d\text{-dim.\ diamond }(D=d+1): & \pbar=\dfrac{(d-1)^2}{2d(d+1)}.
\end{array}
\end{equation}
In particular, 4D diamond has $\pbar=9/40$. Brillouin-zone averages of the walk's own eigenprojectors
give $0.16650$ ($12^3$ grid) and $0.16660$ ($16^3$) for diamond, and $0.22490$ ($8^4$) for 4D diamond.

\emph{Tightness elsewhere.} Comparing with the exact values of Sec.~\ref{sec:model}:
square $0.1673438$ vs.\ $1/6$; triangular $0.2679557$ vs.\ $4/15$;
simple cubic $0.2667272$ vs.\ $4/15$; BCC $0.3215048$ vs.\ $9/28$;
FCC $0.3789146$ vs.\ $25/66$. The bound is within $10^{-3}$ in every case we examined,
and within $10^{-4}$ for simple cubic and BCC.

%=====================================================================
\section{Higher-dimensional trapping integrals}
\label{sec:higherd}
%=====================================================================

\subsection{Moving-shift Grover walk on $\mathbb{Z}^d$}

Consider the Grover coin $G=\tfrac1dJ-I$ on the $2d$ coin states $\pm e_i$ with the moving
shift. In momentum space $S(k)G$ has eigenvalues $\pm1$ for every $k$. Solving
$G\psi=\lambda S(k)^*\psi$ gives flat-band components $(1\pm i\tan\tfrac{k_i}{2})/2$ for $\lambda=+1$
and $(1\mp i\cot\tfrac{k_i}{2})/2$ for $\lambda=-1$. When $k_i$ is uniform on the circle,
$\tan(k_i/2)$ and $\cot(k_i/2)$ are standard Cauchy variables $t_i$. Hence
\begin{equation}
\overline M=\tfrac A2J+\tfrac B2\,\mathrm{blockdiag}\begin{pmatrix}1&-1\\-1&1\end{pmatrix},
\end{equation}
with
\begin{equation}
A_d=\E\Bigl[\frac1{d+\sum_it_i^2}\Bigr]=\int_0^\infty\erfc(\sqrt s)^d\,ds,
\qquad B_d=\frac{1-dA_d}{d}.
\label{eq:Ad}
\end{equation}
The second form of $A_d$ follows from $1/x=\int_0^\infty e^{-sx}ds$ and
$\E[e^{-st^2}]=e^s\erfc\sqrt s$. With $u=(1,\dots,1)/\sqrt{2d}$ and $P_{\rm anti}$ the projector
onto $\{e_{+i}-e_{-i}\}$,
\begin{equation}
\pbar(\phi)=2\bigl[(dA_d)^2|\langle u|\phi\rangle|^2+B_d^2\|P_{\rm anti}\phi\|^2\bigr],
\end{equation}
so that $\pbar(+e_1)=dA_d^2+B_d^2$ and $\pbar(u)=2d^2A_d^2$. The exact values are
\begin{align}
A_2&=\tfrac12-\tfrac1\pi,\\
A_3&=\tfrac12-\tfrac{3-\sq}{\pi}=0.09639923687042003490\ldots,\\
A_4&=\tfrac12-\tfrac{18-8\sq}{3\pi}=0.06035107068870143614\ldots,\\
A_5&=\tfrac12-\tfrac{10}\pi+\tfrac{40\sq}{3\pi}-\tfrac{20\sq\arccos(1/4)}{\pi^2}\nonumber\\
&=0.04153865644854222924\ldots
\end{align}
The $d=2$ case reproduces Inui \emph{et al.}~\cite{Inui2004}. The 3D headline values are
$\pbar(+e_1)=0.08401620467429515245$ and $\pbar(u)=0.1672706316455882907$; in 5D,
$\pbar(+e_1)=0.03373729729788322185$. Because $A_d=\tfrac12\E[\min_kZ_k^2]$ for independent standard
normals, min/max inclusion--exclusion turns $A_5$ into
$\E[\max_{k\le5}Z_k^2]=1+20\sq/\pi-40\sq\arccos(\tfrac14)/\pi^2$. That is a mean-square maximum of
half-normals, the setting considered by Finch~\cite{Finch2011}, so the integral itself is probably
known. $A_5$ closes because it is a moment over a region of the even-dimensional
sphere $S^4$. For $d=6$ ($S^5$) no closed form was found.

\subsection{Weighted coins: a general-$d$ theorem}

Replace the Grover coin by $2|s\rangle\langle s|-I$ with
$|s\rangle=(\sqrt{p_1/2},\sqrt{p_1/2},\dots,\sqrt{p_d/2},\sqrt{p_d/2})$ and $\sum p_i=1$. Let $Z_i$ be
independent standard normals and define
\begin{equation}
A=\tfrac12\E\bigl[\min_kZ_k^2/p_k\bigr]=\int_0^\infty\prod_k\erfc\bigl(\sqrt{p_ku}\bigr)\,du,
\end{equation}
together with $\pi_i=P(\text{coordinate $i$ attains the minimum})$.

\emph{Theorem 3.} Starting in coin state $+e_1$,
\begin{equation}
\pbar=2\Bigl[\tfrac{\pi_1^2}4+\bigl(p_1A-\tfrac{\pi_1}2\bigr)^2+\tfrac{p_1(1-p_1)A^2}2\Bigr].
\end{equation}
The proof uses the Cauchy identities $\E[1/(1+\sum p t^2)]=A$ and
$\E[t_i^2/(1+\sum pt^2)]=\pi_i/p_i-A$.

In $d=3$ everything is elementary. For $\{i,j,k\}=\{1,2,3\}$,
\begin{align}
\pi_i&=\tfrac2\pi\Bigl[\arctan\sqrt{\tfrac{p_i}{p_j}}+\arctan\sqrt{\tfrac{p_i}{p_k}}
-\arctan\sqrt{\tfrac{p_i}{p_jp_k}}\Bigr],\\
A&=\tfrac12\sum_i\Bigl[\tfrac{\pi_i}{p_i}-\tfrac2\pi\Bigl(\tfrac{\sqrt{p_k/p_i}}{1+\sqrt{p_j}}
+\tfrac{\sqrt{p_j/p_i}}{1+\sqrt{p_k}}\Bigr)\Bigr].
\end{align}
These are obtained from solid angles of the three cones meeting at $(\sqrt{p_1},\sqrt{p_2},\sqrt{p_3})$
together with Stokes' theorem for $\int_Tx_i^2\,d\sigma$. For example, $p=(0.2,0.3,0.5)$ gives
$\pbar=0.05362025915918515966$. The $d=2$ case may overlap earlier work on generalized Grover coins.

\subsection{Lackadaisical moving-shift walk}

Add a self-loop of weight $l$ to the moving-shift walk, so that
$|s\rangle=(1,\dots,1,\sqrt l)/\sqrt{2d+l}$~\cite{Wong2015}. The $-1$ flat band disappears. With
\begin{equation}
a(l)=\tfrac12\int_0^\infty e^{-ls/2}\erfc(\sqrt s)^d\,ds,\qquad b=\frac{1-(2d+l)a}{2d},
\end{equation}
the return probability is $l(2d+l)a^2$ from the loop state and $2da^2+2b^2+la^2$ from a moving
state. The closed forms are
\begin{align}
d=2:&\quad a=\tfrac1l\Bigl[1-\tfrac4\pi\tfrac{\arctan q}{q}\Bigr],\qquad q=\sqrt{1+l/2},\\
d=3:&\quad a=\tfrac1l\Bigl[1-3\sqrt{\tfrac{2}{l+2}}\Bigl(1+\tfrac2\pi\Bigl(\arcsin\tfrac2{l+4}\nonumber\\
&\qquad\qquad-2\arcsin\sqrt{\tfrac2{l+4}}\Bigr)\Bigr)\Bigr].
\end{align}
The $d=3$ case rests on the lemma
\begin{align}
\int_0^\infty e^{-bx^2}\erfc(x)^2dx
&=\tfrac12\sqrt{\tfrac\pi b}\Bigl[1+\tfrac2\pi\Bigl(\arcsin\tfrac1{b+1}\nonumber\\
&\qquad-2\arcsin\tfrac1{\sqrt{b+1}}\Bigr)\Bigr],
\end{align}
which is proved via a trivariate Gaussian orthant probability.

\subsection{Moving-shift triangular walk (partial)}

On the triangular lattice with the moving shift, the $-1$ channel from $+e_1$ is
$\tfrac5{12}-\tfrac{4\sq}{3\pi}+\tfrac{57}{16\pi^2}=0.04251818949677224475$. The $+1$ channel is
$0.04312949681537029391$, for which no closed form was found in bases built from
$\pi$, $\sq$, logarithms, Clausen values, Catalan's constant and $\arccos(1/3)$. The total is
$\pbar=0.08564768631214253866$.

%=====================================================================
\section{UST degree laws}
\label{sec:ust}
%=====================================================================

Let $Y$ be the $z\times z$ transfer-current matrix of the edges at a vertex. By the transfer-current
theorem of Burton and Pemantle~\cite{BurtonPemantle1993} and multilinearity of the determinant,
the degree of the vertex in the uniform spanning tree satisfies
\begin{equation}
\sum_kP(\deg=k)\,x^k=\det\bigl[I+(x-1)Y\bigr].
\label{eq:ust-gf}
\end{equation}
The method reproduces the square-lattice values of Manna, Dhar and Majumdar~\cite{Manna1992}.

\emph{Theorem 4 (binomial law).} Suppose the lattice is edge-transitive and the vertex symmetry
acts 2-transitively on its $z$ edges (honeycomb, diamond, Laves graph, $d$-dimensional
diamond). Then
\begin{equation}
\deg\ \stackrel{d}{=}\ 1+\mathrm{Binomial}\Bigl(z-1,\frac1{z-1}\Bigr).
\end{equation}
\emph{Proof.} 2-transitivity forces $Y=aI+b(J-I)$. Foster's theorem gives $a=2/z$, and
Eq.~\eqref{eq:rowsum} gives $a+(z-1)b=1$. So $Y$ has eigenvalue $1$ on the all-ones vector and
$a-b=1/(z-1)$ with multiplicity $z-1$. Hence
$\det[I+(x-1)Y]=x\bigl(\tfrac{z-2}{z-1}+\tfrac{x}{z-1}\bigr)^{z-1}$. $\square$

For the honeycomb this gives $(1/4,1/2,1/4)$ and for diamond $(8/27,4/9,2/9,1/27)$. As
$z\to\infty$ the law tends to $1+\mathrm{Poisson}(1)$, the known mean-field limit.
\emph{The law does not hold} for the pyrochlore, king or snub-square lattices, whose
vertex symmetry is not 2-transitive on the edges. Their distributions are given below.

\emph{Simple cubic.} With $\alpha=(7W+54/(\pi^2W))/36=0.39507915234185137202$,
\begin{align}
P(1)&=8\alpha^3(3\alpha-2)^2,\nonumber\\
P(2)&=-4\alpha^2(3\alpha-2)(30\alpha^2-25\alpha+6),\nonumber\\
P(3)&=2\alpha(360\alpha^4-600\alpha^3+379\alpha^2-108\alpha+12),\nonumber\\
P(4)&=-(2\alpha-1)(360\alpha^4-480\alpha^3+241\alpha^2-52\alpha+4),\nonumber\\
P(5)&=2(2\alpha-1)^2(3\alpha-1)(15\alpha^2-10\alpha+2),\nonumber\\
P(6)&=-(2\alpha-1)^3(3\alpha-1)^2 .
\end{align}
Numerically these are $0.327494821326$, $0.409830450502$, $0.204859445598$, $0.0511275258041$,
$0.00637070295715$ and $0.000317053812167$.

\emph{BCC and FCC.} With $q=1/(\pi^2W)$ for the respective Watson constant,
\begin{align}
P_{\rm BCC}(1)&=\tfrac{3q}{16}(W+4q)^3(4-W-12q)^3\nonumber\\
&=0.339633990106590943,\nonumber\\
P_{\rm BCC}(8)&=\tfrac1{64}(12q-1)(W+4q-2)^3(W+12q-2)^3\nonumber\\
&=1.17884249933455210\times10^{-6},\\
P_{\rm FCC}(1)&=\tfrac{1}{5038848}(5-4W-6q)^2(7-8W-3q)^3\nonumber\\
&\quad\times(8W+3q+1)^3(16W+15q-5)^3\nonumber\\
&=0.349518246140286417 .
\end{align}
The intermediate degrees are polynomials in $W$ and $q$.

\emph{Pyrochlore} ($W=W_f$, $X=1/(\pi^2W)$):
\begin{align}
P(1)&=\tfrac{1}{14155776}[(31-8W-3X)(67+24W+9X)]^2\nonumber\\
&=0.321497700833733694,\nonumber\\
P(6)&=\tfrac{1}{28311552}[(8W+3X+1)(24W+9X-29)]^2\nonumber\\
&=0.000257114517586699387 .
\end{align}

\emph{King lattice} ($p=1/\pi$, $L=\ln(2+\sq)/(\pi\sq)$):
\begin{align}
P(8)&=\tfrac1{128}(1+2p-7L)(4p-1-2L)(6L+12p-5)\nonumber\\
&\quad\times(1+20L-12L^2-32p+48p^2)^2\nonumber\\
&=8.81169962728673452\times10^{-7}.
\end{align}

\emph{Snub square} ($p=1/\pi$), all degrees:
\begin{align}
P(1)&=\tfrac{5}{3888}\bigl[8953-5058\sq+(6336\sq-10692)p\nonumber\\
&\qquad+(1836-1296\sq)p^2\bigr],\nonumber\\
P(2)&=\tfrac1{972}\bigl[30245\sq-51893+(52758-30708\sq)p\nonumber\\
&\qquad+(5940\sq-9666)p^2\bigr],\nonumber\\
P(3)&=\tfrac1{648}\bigl[56447-32500\sq+(28872\sq-50112)p\nonumber\\
&\qquad+(9504-5400\sq)p^2\bigr],\nonumber\\
P(4)&=\tfrac1{972}\bigl[32055\sq-55511+(45522-26172\sq)p\nonumber\\
&\qquad+(4860\sq-8694)p^2\bigr],\nonumber\\
P(5)&=\tfrac1{3888}\bigl[50057-28910\sq+(22608\sq-38988)p\nonumber\\
&\qquad+(7236-4320\sq)p^2\bigr].
\end{align}

\emph{Further two-dimensional laws.} Again $c=\sqrt2\arccos(1/3)/\pi$ and $r=\sq/\pi$.
The 4.8.8 lattice has $P(1)=P(3)=3c(4-3c)/16$ and $P(2)=(9c^2-12c+8)/8$.
The star lattice has $(21/100,29/50,21/100)$.
The kagome lattice has
$P=\tfrac{(5+3r)(17-6r)}{324},\tfrac{18r^2-21r+158}{324},\tfrac{18r^2-21r+77}{324},\tfrac{(4-3r)(1+6r)}{324}$,
with $P(1)-P(4)=P(2)-P(3)=\tfrac14$ exactly.
The dice rim site and the triakis triangle centre are themselves binomial:
$1+\mathrm{Bin}(2,\tfrac14)$ and $1+\mathrm{Bin}(2,\tfrac1{10})$.
For the triangular lattice, $P(1)=(3r-1)(6r-5)^2(6r+1)^2/108$, consistent with the known triangular
sandpile height-one probability~\cite{PonceletRuelle2017}. The general determinantal approach of
Rapoport~\cite{Rapoport2023} covers such degree statistics in principle.
In every case the distribution sums to $1$ and has mean exactly $2$, and $P(1)$ equals $z$
times the sandpile height-one probability (the companion paper).

%=====================================================================
\section{Numerical and symbolic verification methodology}
\label{sec:methods}
%=====================================================================

\emph{Exact Bloch engine.} For a two-dimensional periodic graph with Bloch Laplacian $L(k)$, we
evaluate every needed average $\E_k[\,\cdot\,]$ by an exact residue sum in $z_2=e^{ik_2}$ (the roots
of $\det L$, polished by Newton iteration, with a contour integral around clustered roots near the
node at $k=0$), followed by tanh--sinh quadrature in $k_1$ at 25--60 digits. The three-dimensional
engine integrates $k_3$ exactly and the remaining two dimensions by quadrature. It was validated
on the simple-cubic Watson integral to 12 digits, and BCC Green functions were computed by accelerated
exact binomial series to 45 digits.

\emph{Integer-relation identification.} Closed forms were found with PSLQ~\cite{PSLQ1999} (mpmath
implementation) over small bases. Relations were accepted only when the coefficients were small and
the relation was stable under an increase of precision. Representative precisions: $A_3,A_4$ at 60
digits; $A_5$ at 80 digits (residual $<10^{-80}$); simple-cubic currents at 40 digits (rechecked at
70); FCC at 32 and 50 digits; kagome and triangular at 35--45 digits; the line graph of FCC at 34 digits,
with the energy-8 relations re-verified at 60 digits. In that last case one 34-digit relation proved
\emph{spurious} and was corrected. This cautions against accepting PSLQ output at marginal precision.

\emph{Independent checks.} (a) Symbolic assembly of each channel was compared with the engine's direct
value (typically 20--40 digits). (b) Brillouin-zone diagonalization of the full walk operator $U(k)$ on
grids ($64^3$ for the moving shift, $36^3$ for $L$(SC), $90\times90$ for kagome), averaging each flat
projector. (c) Direct simulation of the walk. Examples: square $0.16734263$ vs.\ $0.16734378$;
triangular $0.2679596$ vs.\ $0.2679557$; honeycomb $0.083337$ vs.\ $1/12$; lazy triangular
$\varepsilon=0.2$ $0.316536$ vs.\ $0.316531$. 3D simulations are memory-limited to $t\lesssim90$, and
their residual decays like $t^{-3/2}$, as expected for the escaping part of the wave. (d) Closed-form
Green functions were checked against the lattice equations, for example
$6G'(x)+\tfrac12\sum_{\rm nbrs}G'(x+e)=\delta_{x,0}$ on FCC with residual $3\times10^{-42}$.
(e) Every UST distribution was checked to sum to 1 with mean exactly 2 (to $10^{-16}$), and against
the independent sandpile determinant.

%=====================================================================
\section{Discussion and conclusion}
\label{sec:conclusion}
%=====================================================================

The transfer-current principle, Eq.~\eqref{eq:principle}, turns Grover-walk localization into a
problem of classical electrical networks. This explains why the trapping probabilities involve exactly
the constants of lattice potential theory: Watson integrals, Joyce's elliptic forms, and Horiguchi's
triangular Green function. It also explains why the signless Laplacian, that is, the Green function at the
\emph{opposite} band edge, governs the $-1$ channel of non-bipartite lattices. A by-product is the
far-edge FCC identity, Eq.~\eqref{eq:fcc-far}, and the real-parameter triangular Green function,
Eq.~\eqref{eq:tri-real}.

The universal bound, Theorem~2, is sharp on 2-arc-transitive bipartite lattices and within $10^{-3}$ of
every exact value we computed. So localization in flip-flop Grover walks is controlled, to excellent
accuracy, by the vertex degree alone. The same symmetry yields the binomial UST degree law of
Theorem~4, which connects quantum-walk localization to spanning-tree combinatorics.

Open problems include a closed form for the $+1$ channel of the moving-shift triangular walk, the
integral $A_6$, and the $-1$ channel of the anisotropic triangular walk. The latter needs a signless
Green function whose elliptic parameter depends on the weights.

\begin{acknowledgments}
The computations and the first draft of this manuscript were carried out with the assistance
of an AI coding agent (Claude Code). All numerical code is available in the accompanying
repository.
\end{acknowledgments}

\begin{thebibliography}{99}

\bibitem{Kempe2003}
J.~Kempe,
Quantum random walks: An introductory overview,
\href{https://doi.org/10.1080/00107151031000110776}{Contemp.\ Phys.\ \textbf{44}, 307 (2003)},
arXiv:quant-ph/0303081.

\bibitem{Szegedy2004}
M.~Szegedy,
Quantum speed-up of Markov chain based algorithms,
in \emph{Proceedings of the 45th Annual IEEE Symposium on Foundations of Computer Science (FOCS 2004)}
(IEEE, 2004), pp.~32--41,
\href{https://doi.org/10.1109/FOCS.2004.53}{doi:10.1109/FOCS.2004.53}.

\bibitem{Wong2015}
T.~G.~Wong,
Grover search with lackadaisical quantum walks,
\href{https://doi.org/10.1088/1751-8113/48/43/435304}{J.\ Phys.\ A: Math.\ Theor.\ \textbf{48}, 435304 (2015)}.

\bibitem{Inui2004}
N.~Inui, Y.~Konishi, and N.~Konno,
Localization of two-dimensional quantum walks,
\href{https://doi.org/10.1103/PhysRevA.69.052323}{Phys.\ Rev.\ A \textbf{69}, 052323 (2004)}.

\bibitem{Higuchi2014}
Yu.~Higuchi, N.~Konno, I.~Sato, and E.~Segawa,
Spectral and asymptotic properties of Grover walks on crystal lattices,
\href{https://doi.org/10.1016/j.jfa.2014.09.003}{J.\ Funct.\ Anal.\ \textbf{267}, 4197 (2014)}, arXiv:1401.0154.

\bibitem{Kirchhoff1847}
G.~Kirchhoff,
\"Uber die Aufl\"osung der Gleichungen, auf welche man bei der Untersuchung der linearen
Vertheilung galvanischer Str\"ome gef\"uhrt wird,
\href{https://doi.org/10.1002/andp.18471481202}{Ann.\ Phys.\ Chem.\ \textbf{72}, 497 (1847)}.

\bibitem{BurtonPemantle1993}
R.~Burton and R.~Pemantle,
Local characteristics, entropy and limit theorems for spanning trees and domino tilings via
transfer-impedances,
\href{https://doi.org/10.1214/aop/1176989121}{Ann.\ Probab.\ \textbf{21}, 1329 (1993)}.

\bibitem{Watson1939}
G.~N.~Watson,
Three triple integrals,
\href{https://doi.org/10.1093/qmath/os-10.1.266}{Q.\ J.\ Math.\ \textbf{os-10}, 266 (1939)}.

\bibitem{GlasserZucker1977}
M.~L.~Glasser and I.~J.~Zucker,
Extended Watson integrals for the cubic lattices,
\href{https://doi.org/10.1073/pnas.74.5.1800}{Proc.\ Natl.\ Acad.\ Sci.\ USA \textbf{74}, 1800 (1977)}.

\bibitem{Joyce1994}
G.~S.~Joyce,
On the cubic lattice Green functions,
\href{https://doi.org/10.1098/rspa.1994.0072}{Proc.\ R.\ Soc.\ Lond.\ A \textbf{445}, 463 (1994)}.

\bibitem{Joyce1998}
G.~S.~Joyce,
On the cubic modular transformation and the cubic lattice Green functions,
J.\ Phys.\ A: Math.\ Gen.\ \textbf{31}, 5105 (1998).

\bibitem{Horiguchi1972}
T.~Horiguchi,
Lattice Green's functions for the triangular and honeycomb lattices,
\href{https://doi.org/10.1063/1.1666155}{J.\ Math.\ Phys.\ \textbf{13}, 1411 (1972)}.

\bibitem{Foster1949}
R.~M.~Foster,
The average impedance of an electrical network,
in \emph{Reissner Anniversary Volume: Contributions to Applied Mechanics}
(J.~W.~Edwards, Ann Arbor, 1949), pp.~333--340.

\bibitem{Finch2011}
S.~R.~Finch,
Mean width of a regular cross-polytope,
arXiv:1112.0499 (2011).

\bibitem{Kenyon2002}
R.~Kenyon,
The Laplacian and Dirac operators on critical planar graphs,
\href{https://doi.org/10.1007/s00222-002-0249-4}{Invent.\ Math.\ \textbf{150}, 409 (2002)}.

\bibitem{Manna1992}
S.~S.~Manna, D.~Dhar, and S.~N.~Majumdar,
Spanning trees in two dimensions,
\href{https://doi.org/10.1103/PhysRevA.46.R4471}{Phys.\ Rev.\ A \textbf{46}, R4471 (1992)}.

\bibitem{PonceletRuelle2017}
A.~Poncelet and P.~Ruelle,
Sandpile probabilities on triangular and hexagonal lattices,
\href{https://doi.org/10.1088/1751-8121/aa9255}{J.\ Phys.\ A: Math.\ Theor.\ \textbf{51}, 015002 (2018)},
arXiv:1708.03493.

\bibitem{Rapoport2023}
A.~Rapoport,
Correlations in uniform spanning trees: a fermionic approach,
arXiv:2312.14992 (2023).

\bibitem{PSLQ1999}
H.~R.~P.~Ferguson, D.~H.~Bailey, and S.~Arno,
Analysis of PSLQ, an integer relation finding algorithm,
\href{https://doi.org/10.1090/S0025-5718-99-00995-3}{Math.\ Comp.\ \textbf{68}, 351 (1999)}.

\end{thebibliography}

\end{document}
